\section{Genuine trace growth in Yang--Mills theory}
\label{sec:trace-hierarchy}

This section derives the expectation-value evolution summarized
in the exposition. The contour is the actual Loewner prefix, not
a reference contour deformed by an inverse map. We fix the theory
and an observable before differentiating. The results are exact
smooth-connection identities and their conditional quantum
averages; no closed-form path integral is required.

\subsection{Theory, field resolution and averaging assumptions}

Fix Euclidean $SU(N)$ Yang--Mills, $N\ge2$, at $\theta=0$.
With Hermitian generators normalized by
$\tr(T^aT^b)=\delta^{ab}/2$, use
\begin{align}\label{eq:trace-action}
 F_{\mu\nu}[A]&=\partial_\mu A_\nu-\partial_\nu A_\mu
                         -\ii[A_\mu,A_\nu],\\
 D_\mu^A X&=\partial_\mu X-\ii[A_\mu,X],\qquad
 S[A]=\frac1{2g^2}\int_{\R^4}\tr F_{\mu\nu}[A]F_{\mu\nu}[A]\dd^4x.
 \notag
\end{align}
The action, coupling and quantum state stay fixed. Embed the
Loewner upper half-plane in $x_3=x_4=0$; removing its hull does
not change the physical spacetime or the measure.

To define a smooth probe, take the gauge-covariant field flow
\begin{equation}\label{eq:trace-field-flow}
 \partial_\tau B_\mu=D_\nu^B F_{\nu\mu}[B],\qquad B_0=A,
 \qquad \tau>0\ \hbox{fixed}.
\end{equation}
The Wilson observable is constructed from $B_\tau[A]$ and averaged
in the original interacting measure. There is no independent
Gaussian measure for $B_\tau$. Field-flow time $\tau$, Loewner
capacity time $t$ and the arithmetic shift $\omega$ are distinct.
Positive-flow local observables and their perturbative
renormalization are discussed in
\cite{LuscherFlow,LuscherWeiszFlow}.

\begin{assumption}[Regularity for the expectation identities]
\label{ass:trace-regularity}
The chosen regulated or flowed construction defines the indicated
holonomies and insertion correlators on the swept compact region.
Connections are sufficiently smooth for ordered-transport
differentiation; expectations of the required local field bounds
justify passing derivatives through the expectation. For the
short-time remainder we require the two finite moments specified
in Proposition~\ref{prop:trace-remainder}.
\end{assumption}

This assumption is not a construction of nonperturbative
continuum four-dimensional Yang--Mills. Results on local
positive-flow observables alone do not establish it for every
cusped nonlocal loop. Without positive flow the same algebra
gives bare regulated identities, whose passage to renormalized
unflowed observables must include perimeter, cusp and contact
terms. We do not perform that passage here. Below we suppress
$\tau$ and write $A,F,D$ for the selected smooth connection;
this does not identify $B_\tau$ with the integration variable in
a Schwinger--Dyson equation.

\subsection{The actual prefix and its return chord}

Use half-plane capacity $2t$ and a deterministic driver $u(t)$:
\begin{equation}\label{eq:trace-definition}
 \partial_tg_t(z)=\frac2{g_t(z)-u(t)},\qquad
 q(t)=\lim_{y\downarrow0}g_t^{-1}(u(t)+\ii y),\qquad q(0)=0.
\end{equation}
We assume a trace that is $C^1$ for positive $t$, with the local
integrability needed at zero. Section~\ref{subsec:linear-trace}
supplies an explicit analytic example. No Brownian trace or
contour average is assumed; general driver regularity is a
separate issue~\cite{LindTran}.

Let $U_t$ transport along $q[0,t]$, with later parameters ordered
to the left. For $0\le r\le1$ let $R_t(r)$ transport from $0$ to
$rq(t)$ along the radial segment, and put
\begin{equation}\label{eq:trace-holonomy}
 \mathsf H_t=R_t(1)^{-1}U_t,\quad
 W_\tau(t)=\frac1N\vev{\tr\mathsf H_t},\quad
 \widehat F_t(r)=R_t(r)^{-1}F_{12}(rq(t))R_t(r).
\end{equation}
The straight return chord closes the trace's color indices.
It is a choice of observable, not an identity for every possible
endpoint state. The trace is in the fundamental representation;
basepoint conjugation leaves all the displayed traced insertions
gauge invariant. At zero size $\mathsf H_0=I$ and $W_\tau(0)=1$.

Pure gauge transport reverses to its inverse. The same assertion
does not follow for a scalar-arclength-coupled $\mathcal N=4$
Wilson line: its scalar term requires an additional analysis.
The chord-completed observable therefore supplies a pure-YM
calculation within a fixed theory, not an implicit supersymmetric
extension or a proposed all-input arithmetic transfer.

\subsection{Exact cancellation of the moving-tip term}

Define
\begin{equation}\label{eq:trace-J}
 k_t=q_1\dot q_2-q_2\dot q_1,\qquad
 \mathsf J_t(r)=\int_0^r s\widehat F_t(s)\dd s,\qquad
 \mathsf J_t=\mathsf J_t(1).
\end{equation}

\begin{proposition}[Growing-trace transport]
\label{prop:trace-transport}
For a smooth Hermitian connection and the specified completion,
\begin{equation}\label{eq:trace-matrix-flow}
 \dot{\mathsf H}_t=\ii k_t\mathsf J_t\mathsf H_t.
\end{equation}
Under Assumption~\ref{ass:trace-regularity},
\begin{equation}\label{eq:trace-first-expectation}
 \dot W_\tau=\ii k_t M_1,\qquad
 M_1=\frac1N\vev{\tr(\mathsf J_t\mathsf H_t)}.
\end{equation}
\end{proposition}
\begin{proof}
Only the newly added prefix segment contributes to
\[
 \dot U_t=\ii A_\mu(q(t))\dot q^\mu U_t.
\]
For the outward chord $x(s)=s q(t)$, its displacement is
$s\dot q(t)$. The gauge part of the variation formula
\eqref{eq:variation} has
$F_{\nu\mu}\dot q^\nu q^\mu=-k_tF_{12}$, and hence
\begin{equation}\label{eq:trace-radial-variation}
 R_t(r)^{-1}\dot R_t(r)
 =\ii rR_t(r)^{-1}A_\mu(rq)\dot q^\mu R_t(r)
                                  -\ii k_t\mathsf J_t(r).
\end{equation}
At $r=1$ the endpoint connection cancels exactly when
$R_t(1)^{-1}U_t$ is differentiated. The remaining term is
\eqref{eq:trace-matrix-flow}. Tracing and averaging in the fixed
measure gives \eqref{eq:trace-first-expectation}. No field
equation, Gaussian contraction or factorization has been used.
\end{proof}

For a planar curve with this chord, its oriented area satisfies
\begin{equation}\label{eq:trace-area}
 \mathcal A(t)=\frac12\int_0^t k_v\dd v.
\end{equation}
A constant commuting curvature gives
$\mathsf H_t=\exp(\ii F_{12}\mathcal A(t))$, checking the
orientation and the factor of two independently.

\subsection{Curvature transport and the next expectation equation}

Differentiating the conjugation in \eqref{eq:trace-holonomy}
with \eqref{eq:trace-radial-variation} gives
\begin{equation}\label{eq:trace-curvature-flow}
 \partial_t\widehat F_t(r)
 =r\dot q^\mu\widehat{D_\mu F_{12}}_t(r)
                      +\ii k_t[\mathsf J_t(r),\widehat F_t(r)].
\end{equation}
The endpoint connection in the conjugation converts the ordinary
derivative of $F$ into its covariant derivative. Define
\begin{align}\label{eq:trace-second-moments}
 \mathsf D_{\mu,t}&=\int_0^1r^2\widehat{D_\mu F_{12}}_t(r)\dd r,
 &\mathsf R_t&=\int_0^1r[\mathsf J_t(r),\widehat F_t(r)]\dd r,\\
 \mathsf P_t&=\int_{0<s<r<1}rs\widehat F_t(s)\widehat F_t(r)\dd s\dd r.
 \notag
\end{align}
In particular $\dot{\mathsf J}_t=\dot q^\mu\mathsf D_{\mu,t}
+\ii k_t\mathsf R_t$. Splitting the square of its defining
integral into the two triangles yields
\begin{equation}\label{eq:trace-ordering}
 \mathsf J_t^2+\mathsf R_t=2\mathsf P_t.
\end{equation}
Indeed, on $s<r$, the square contributes
$rs(\widehat F_t(r)\widehat F_t(s)+\widehat F_t(s)\widehat F_t(r))$,
and the commutator contributes their difference with the second
ordering positive. Thus the smaller radial argument is on the left
in $\mathsf P_t$.

\begin{proposition}[The first two quantum equations]
Under Assumption~\ref{ass:trace-regularity}, let
\[
 M_{D,\mu}=\frac1N\vev{\tr(\mathsf D_{\mu,t}\mathsf H_t)},\qquad
 M_2=\frac1N\vev{\tr(\mathsf P_t\mathsf H_t)}.
\]
Then, with $\mu=1,2$ summed,
\begin{align}\label{eq:trace-hierarchy}
 \dot W_\tau&=\ii k_t M_1, &
 \dot M_1&=\dot q^\mu M_{D,\mu}+2\ii k_t M_2,\\
 \ddot W_\tau&=\ii\dot k_tM_1+\ii k_t\dot q^\mu M_{D,\mu}
                                      -2k_t^2M_2.
 \notag
\end{align}
\end{proposition}
\begin{proof}
The first equation is Proposition~\ref{prop:trace-transport}.
Differentiate $\mathsf J_t\mathsf H_t$ and use
\eqref{eq:trace-ordering}. Its derivative is
$(\dot q^\mu\mathsf D_{\mu,t}+2\ii k_t\mathsf P_t)\mathsf H_t$.
The second equation follows by tracing and averaging. The product
rule gives the third on positive-time intervals with the required
derivatives.
\end{proof}

The equations retain the complete interacting expectation.
The derivative and two-curvature insertions are new unknowns.
Writing $\dot W_\tau/W_\tau$ as an unspecified coefficient is
not a closure of the hierarchy.

\subsection{The explicit trace of a linear driver}
\label{subsec:linear-trace}

Choose $u(t)=at$, $a\in\R$. At fixed initial $z$ put
$w_s=g_s(z)-as$, so $\dot w_s=2/w_s-a$. For $a\ne0$ a primitive
of $w/(2-aw)$ is
\[
 -\frac wa-\frac2{a^2}\log(2-aw).
\]
At the tip $w_0=q(t)$ and $w_t=0$. Integrating on the branch
issuing into the upper half-plane gives
\begin{equation}\label{eq:trace-linear-implicit}
 aq+2\log(1-aq/2)=a^2t,\qquad \dot q=a-2/q.
\end{equation}
The local tip ODE is special to this driver; a general trace
depends on the driver's history. In particular it is not obtained
by substituting the singular tip into the forward map equation.

Set $s=\sqrt t$. The identity $q q_s=2asq-4s$ determines
\begin{equation}\label{eq:trace-linear-series}
 q(t)=2\ii\sqrt t+\frac{2a}{3}t-\frac{\ii a^2}{18}t^{3/2}
                +\frac{a^3}{135}t^2+O_a(t^{5/2}).
\end{equation}
For a direct analyticity argument, put $q=sv$ in
\eqref{eq:trace-linear-implicit}, divide out its double zero in
$s$, and apply the implicit-function theorem at $v=2\ii$;
the derivative in $v$ is nonzero. Therefore the coefficient
recursion describes a local analytic series in $s$, not only a
formal asymptotic ansatz. It follows that
\begin{align}\label{eq:trace-area-series}
 \mathcal A(t)&=-\frac{2a}{9}t^{3/2}
                    -\frac{11a^3}{1350}t^{5/2}+O_a(t^{7/2}),\\
 k_t&=-\frac{2a}{3}\sqrt t
                    -\frac{11a^3}{270}t^{3/2}+O_a(t^{5/2}).
 \notag
\end{align}
The negative sign for $a>0$ records clockwise orientation. For
$a=0$, $q(t)=2\ii\sqrt t$ and the return chord exactly retraces
the prefix. Thus $\mathsf H_t=I$ for every field configuration.
This is a control for the selected observable, not a universal
statement about straight-line observables.

\subsection{A quantitative short-time expectation bound}

For a fixed small terminal time $t$, let $K_0$ bound $\norm{F_{12}}$
on all swept chords, and let $K_1$ bound
$\norm{e^\mu D_\mu F_{12}}$ for every unit planar vector $e$ on
that region. These bounds may depend on the field configuration.
Put
\begin{equation}\label{eq:trace-geometric-integrals}
 I(t)=\int_0^t|k_v|\dd v,\qquad
 I_1(t)=\int_0^t|k_v|\,|q(v)|\dd v,
 \qquad F_0=F_{12}(0).
\end{equation}

\begin{proposition}[Remainder with explicit constants]
\label{prop:trace-remainder}
For each smooth Hermitian connection,
\begin{equation}\label{eq:trace-remainder}
 \left|\frac1N\tr\mathsf H_t-1+
             \frac{\mathcal A(t)^2}{2N}\tr F_0^2\right|
 \le\frac{K_0K_1}{6}I(t)I_1(t)+\frac{K_0^3}{48}I(t)^3.
\end{equation}
If $\vev{K_0K_1}$ and $\vev{K_0^3}$ are finite on a fixed small
swept neighborhood, the averaged inequality holds with these
expected moments on the right. In particular, at fixed positive
field-flow time and for the driver $u(t)=at$,
\begin{equation}\label{eq:trace-quantum-expansion}
 W_\tau(t)=1-\frac{2a^2}{81}C_\tau t^3+O_{\tau,a}(t^{7/2}),
 \qquad C_\tau=\frac1N\vev{\tr F_{12}[B_\tau](0)^2}.
\end{equation}
\end{proposition}
\begin{proof}
Radial differentiation of the conjugated curvature gives
\[
 \norm{\widehat F_v(r)-F_0}\le r|q(v)|K_1,\qquad
 \norm{\mathsf J_v-F_0/2}\le |q(v)|K_1/3.
\]
Set $\mathsf M=\int_0^t k_v\mathsf J_v\dd v$ and
$S=K_0I(t)/2$. Then
\[
 \norm{\mathsf M-\mathcal A(t)F_0}\le K_1I_1(t)/3,
 \qquad \norm{\mathsf M}\le S,\qquad |\mathcal A(t)|\le I(t)/2.
\]
Each $\mathsf J_v$ is traceless, so the linear Dyson coefficient
has zero trace. By cyclicity, the trace of the second ordered
coefficient is $-\tr\mathsf M^2/2$. The anti-Hermitian generator
$\ii k_v\mathsf J_v$ gives unitary evolution. The third-order
integral remainder, with the remaining exact propagator retained,
therefore has norm at most $S^3/6$; no exponential moment is needed.
Finally,
\begin{align*}
 \frac1N\left|\tr(\mathsf M^2-\mathcal A^2F_0^2)\right|
 &\le\norm{\mathsf M-\mathcal A F_0}
                   (\norm{\mathsf M}+|\mathcal A|K_0)\\
 &\le K_0K_1I I_1/3.
\end{align*}
Dividing this contribution by two and adding $S^3/6$ proves
\eqref{eq:trace-remainder}. Averaging uses the stated moments.
For the linear driver, $I=O_a(t^{3/2})$ and $I_1=O_a(t^2)$;
\eqref{eq:trace-area-series} supplies the coefficient in
\eqref{eq:trace-quantum-expansion}.
\end{proof}

No Wick expansion evaluates $C_\tau$. It is a local moment of
the same interacting theory and is nonnegative in a positive
Euclidean measure. The order of limits is essential: $t\downarrow0$
with $\tau$ fixed. Removing the resolution first can invalidate
the moment bounds and introduces unflowed cusp and short-distance
effects. This proposition is not a uniform smoothing-removal
theorem or an unsmeared-loop OPE theorem.

\subsection{Schwinger--Dyson equations and the unresolved closure}

At a finite lattice regulator, product Haar measure and the
Wilson action
$S_{\rm lat}=\beta\sum_p(1-\Ree\tr U_p/N)$ give the exact
integration-by-parts identity
\begin{equation}\label{eq:trace-lattice-sd}
 \vev{\nabla_e^a f}=\vev{f\nabla_e^aS_{\rm lat}},
\end{equation}
where $\nabla_e^a$ is a left-invariant link derivative. Applying
it to Wilson observables gives deformation and splitting/joining
terms. Finite-$N$ expectations of products remain part of the
system. Factorizing them is an additional large-$N$ reduction,
not a prerequisite for the quantum identity itself.

The continuum formal loop equation relates a covariant divergence
of an area insertion to contact terms; its loop-space Laplacian
form and supersymmetric extension are discussed in
\cite[Lecture~3]{MakeenkoLoops}. Neither operator is automatically
the first derivative along the specified Loewner family.
The first difference can be exhibited at the field-component
level. In unflowed notation set $E_\nu=\sum_{\mu=1}^4D_\mu F_{\mu\nu}$.
Then
\begin{align}\label{eq:trace-transverse-residual}
 \dot q^1D_1F_{12}+\dot q^2D_2F_{12}
 ={}&\dot q^1E_2-\dot q^2E_1\\
 &-\dot q^1(D_3F_{32}+D_4F_{42})
   +\dot q^2(D_3F_{31}+D_4F_{41}).\notag
\end{align}
Consequently an equation for $E_\nu$ leaves an explicit transverse
insertion in $\dot q^\mu M_{D,\mu}$: the second line, radially
transported, integrated with weight $r^2$, followed by
$\mathsf H_t$, traced and averaged. Planar placement of the
contour does not delete it.

For a local algebraic control, take $T_3=\sigma_3/2$ in $SU(2)$,
$A_2=b(x_1^2-x_3^2)T_3/2$, and all other components zero. Then
\begin{equation}\label{eq:trace-eom-control}
 F_{12}=b x_1T_3,\quad F_{32}=-b x_3T_3,\quad E_\nu=0,
 \quad D_1F_{12}=bT_3=-D_3F_{32}.
\end{equation}
For $b=1$ the deleted component would have operator norm $1/2$
despite the vanishing total equation-of-motion insertion. This
is a counterexample to a component replacement, not a finite-action
vacuum configuration or a no-go theorem for averaged closure.
Any symmetry that removes its averaged counterpart must be proved
for the chosen observable.

At positive flow time, $E[B_\tau]$ is not the functional derivative
of $S[A]$. Applying \eqref{eq:trace-lattice-sd} differentiates
$B_\tau[A]$ and produces its response kernel. The bare unflowed
contact equation cannot be imported unchanged. The additional
two-curvature moment, transverse insertion, flow-response kernel
and possible unflowed cusp limit have not been reduced to $W_\tau$
and $M_1$ alone.

An enlarged system including loops and open lines has been derived
for lattice Yang--Mills--Higgs~\cite{ShenSmithZhuLines}. That is
methodological precedent with a different action, not a closure
theorem for the fixed pure-YM observable here. The bounded next
question is to evaluate the Schwinger--Dyson reduction of the
specific derivative moment while retaining the displayed residuals.
No closed expectation-value evolution, arithmetic realization or
positive cumulative norm is claimed by the present hierarchy.
